\section{Introduction}\label{sec:introduction}

Let $G$ be a discrete group. Its complex Banach group algebra
$\ell^1(G)$ consists of functions $a\colon G\to\C$ with
$\|a\|_1=\sum_{g\in G}|a(g)|<\infty$, equipped with convolution
\[
 (a*b)(g)=\sum_{h\in G}a(h)b(h^{-1}g).
\]
Write $\mathcal C(G)$ for the set of conjugacy classes, and
$\mathcal C_{\rm fin}(G)$ for those consisting of finite-order elements.
The subscript refers to element order, not to the size of a conjugacy
class.

For an idempotent matrix $e=(e_{ij})\in M_q(\ell^1(G))$, its
Hattori--Stallings trace \cite{Hattori1965,Stallings1965} has coefficients
\begin{equation}\label{intro:trace}
 \tau_C(e)=\sum_{g\in C}\sum_{i=1}^q e_{ii}(g),
 \qquad C\in\mathcal C(G).
\end{equation}
These sums converge absolutely, and
$\sum_C|\tau_C(e)|\le\sum_i\|e_{ii}\|_1$.
The trace is additive under direct sums and depends only on the
isomorphism class of the projective module $e\ell^1(G)^q$.
It therefore defines
\[
 \operatorname{HS}^1\colon K_0(\ell^1(G))\longrightarrow
                           \ell^1(\mathcal C(G)),
\]
where $K_0$ is the Grothendieck group of finitely generated projective
right modules. The $\ell^1$-Bass conjecture
\cite[Conjecture~2.2]{BCM2004} asserts that this map takes values in the
algebraic direct sum over $\mathcal C_{\rm fin}(G)$.
We prove this assertion for all discrete groups.

\begin{theorem}\label{thm:main}
Let $G$ be a discrete group, let $q\ge1$, and let
$e\in M_q(\ell^1(G))$ be idempotent.
There is a finite set $F\subset\mathcal C_{\rm fin}(G)$ such that
$\tau_C(e)=0$ for every $C\notin F$. Equivalently,
\[
 \operatorname{HS}^1\bigl(K_0(\ell^1(G))\bigr)
 \subseteq \bigoplus_{C\in\mathcal C_{\rm fin}(G)}\C[C].
\]
\end{theorem}

The finite-support assertion is part of the conclusion, in addition to
vanishing on infinite-order classes. It is not an automatic consequence
of absolute summability. In particular, the proof must make its
geometric cutoff independent of the conjugacy class under consideration.

\subsection{History and relation to earlier methods}
The Hattori--Stallings trace refines the rank of a projective module by
taking values in the quotient of its ring by additive commutators
\cite{Hattori1965,Stallings1965}. For the algebraic group rings $\Z G$
and $\C G$, this quotient is indexed by conjugacy classes. Bass's
conjecture for the integral group
ring asks that this trace be supported at the identity
\cite{Bass1976}; the complex-group-ring formulation permits
finite-order classes. These support questions connect the algebraic
structure of projective modules with the conjugacy structure of a group.

Cyclic homology gives one route to these questions. Burghelea's
computation separates group-ring cyclic homology into contributions
from finite- and infinite-order conjugacy classes
\cite{Burghelea1985}. Eckmann used this description and the Chern
character to prove Bass-type vanishing for several classes of groups
\cite{Eckmann1986}. Emmanouil extended this method through nilpotence
properties of the cyclic periodicity operator \cite{Emmanouil1998}.
Our argument retains cyclic symmetry but works directly with weighted
configurations and differential forms.

Berrick, Chatterji, and Mislin formulated the $\ell^1$-Bass conjecture
in the finite-support form used here. They proved it for countable
groups whose Bost assembly map is rationally surjective in degree zero
\cite[Theorem~1.4]{BCM2004}. Together with Lafforgue's Banach-algebra
$K$-theory \cite{Lafforgue2002}, this includes amenable groups;
the passage from countable subgroups to arbitrary groups is explained
in \cite[Lemma~1.5]{BCM2004}. Theorem~\ref{thm:main} removes the
assembly hypothesis while retaining the algebraic finite-support
conclusion.

Ji, Ogle, and Ramsey developed related Bass conjectures for
rapid-decay completions and higher Chern characters
\cite{JOR2010,JOR2014}. Their work highlights the growth estimates
needed when passing from group rings to analytic completions.
Here the summability needed in the proof comes from a similarity
reduction to an exponentially decaying idempotent. The subsequent
estimates use a spatial scale uniform over all classes to be excluded.

The weighted cyclic configurations and Pfaffian detection strategy
build on the complex-group-ring argument in
\cite[Sections~7--8]{OpenAIBass2026}. In the group-ring setting,
nonzero products of coefficients involve group elements from a fixed
finite set. For an $\ell^1$ idempotent these elements may have
arbitrarily large word length, even after reduction to exponential
decay. The estimates must retain summable length factors without
losing control as the number of factors increases. We obtain them
from smooth spatial summaries whose derivative bounds are independent
of the size of the input support, multiscale time averaging with a
positive translation component, and a final ordered-kernel bound
uniform in the closing group element. A single spatial radius then
excludes all but finitely many finite-order classes. All ingredients
used here are proved below; the complex-group-ring theorem is not an
input. Standard concentration and transgression arguments are
credited at their points of use.

\subsection{Proof strategy}
The argument uses Banach algebra power series, differential forms on
finite-dimensional simplices, and elementary estimates for integral
operators. We give the needed estimates in full.

First, a finite-support approximation followed by an analytic
idempotent correction replaces $e$ by a similar idempotent $p$ with an
exponential moment in a finitely generated subgroup
(Section~\ref{sec:reduction}). This step preserves every trace
coefficient. It does not assert that $p$ has finite support.

For a fixed $g\in G$, consider cyclic lists of group elements whose
closing edge is twisted by multiplication by $g$. Products of the
coefficients of $p$ assign complex weights to these lists. Their total
weight is $\tau_{[g]}(p)$, and idempotency makes the weights compatible
with deleting vertices. The ordered times of a list form a simplex.
On these simplices, a one-form invariant under simultaneous time
translation and equal to one on the translation vector will be called
a connection form. A cyclic Stokes argument shows that the weighted
average of its curvature Pfaffian equals the trace coefficient, up to
sign (Section~\ref{sec:cyclic}).

The task is then to construct compatible connection forms with small
averaged curvature. Interpret a timed list as a step function with
values in $G$, and average its Dirac masses on a sequence of time scales.
The resulting probability distributions are summarized by smooth
subprobability distributions on auxiliary labels. These summaries
record spatial concentration while controlling derivatives caused by
moving one boundary of the step function
(Section~\ref{sec:channels}). Pairing their differentials with bounded
odd time kernels gives one-forms. Their contractions with time
translation telescope across the averaging scales. A final equality
kernel detects the twist and leaves a strictly positive contraction,
which allows normalization to a connection
(Section~\ref{sec:connection}).

The curvature estimates use three different features of this
construction. At the initial scale, a curvature entry vanishes unless
its two times are close; an occupancy estimate controls the resulting
sparse Pfaffians on average. At intermediate scales, the short support
of the time kernel gives a Hilbert--Schmidt bound. At the last scale,
the equality kernel reduces to weighted order kernels, whose singular
values satisfy $s_\nu\le C/\nu$. Expanding the Pfaffian across these
stages gives exponential decay in the simplex dimension
(Section~\ref{sec:curvature}). A single spatial radius works for every
infinite-order class and every finite-order class avoiding a fixed
finite word ball. This uniformity proves the finite-support conclusion
in Section~\ref{sec:conclusion}.
